The assumed physical model is a coherent single electron with constant effective
mass $m^*>0$ in one spatial dimension:
\[
i\hbar\partial_t\psi(x,t)=-\frac{\hbar^2}{2m^*}\partial_x^2\psi(x,t)+V(x)\psi(x,t).
\]
Here $x$ is position (internally metres), $t$ is time (seconds), $\hbar$ is the
reduced Planck constant (joule seconds), $V$ is potential energy (joules), and
$\psi$ is a position amplitude (metres$^{-1/2}$) for a normalized packet.
The domain has a finite device region and uniform semi-infinite leads. The
public API uses nm, meV and ps with explicit conversion. This equation is an
assumption, not a derivation from classical mechanics.

For elastic propagation, let $t(E)$ and $r(E)$ be complex amplitudes in the
unit-flux incoming/outgoing lead basis at explicitly fixed reference planes.
Then $T(E)=|t(E)|^2$ and $R(E)=|r(E)|^2$ are dimensionless probabilities for the
open channel. Phase is retained in the amplitudes. Let $A(E)$ be an
energy-normalized incoming amplitude (meV$^{-1/2}$ at the API), with
$\int_{\mathcal E}|A(E)|^2\,dE=1$ over the declared finite interval
$\mathcal E$ of propagating lead energies. Define
\[
P=\int_{\mathcal E}|A(E)t(E)|^2\,dE,\qquad
b(E)=\frac{A(E)t(E)}{\sqrt P},\qquad
F=\left|\int_{\mathcal E}b_{\rm target}(E)^*b(E)\,dE\right|^2.
\]
$b$ and $b_{\rm target}$ are separately normalized conditional transmitted
states; $F$ is undefined when either cannot be normalized. $P$ and $F$ must be
reported separately. A global phase cancels from $F$. Optional delay alignment
uses a declared bounded lag search; its bound, sign, resolution and boundary
status accompany raw $F$. Numerical quadrature and finite-window normalization
are approximations with recorded settings.
